\subsection{Related Work}
\label{Section: related work}
\begin{table}[t]
    \centering
    \caption{Comparison of optimization and convergence guarantee of deep hashing methods.}
    \label{Table: summary of deep hashing algorithms}
    \resizebox{\linewidth}{!}{
\begin{tabular}{cccccc}
\toprule
\textbf{Type} & \textbf{Method} & \textbf{Model} 
  & \textbf{Quantization Loss} & \textbf{Optimization} & \textbf{Complexity} \\
\midrule
    \multirow{4}{*}{\centering Discrete} 
    & DDSH \cite{liDeepSupervisedHashing} & $\min L_s(\bb) + Q(\bu, \bb)$ & \centering $\|\cdot\|^2$ & BQP & \multirow{14}{*}{\centering ---} \\
    \cmidrule(lr){2-5}
    & DSDH \cite{li2017deep}& \multirow{3}{*}{\centering $\min L_s(\bu) + Q(\bu, \bb)$} & \multirow{3}{*}{\centering $\|\cdot\|^2$} & $sgn$ update &  \\ 
    & DPSH \cite{liFeatureLearningBased2016} & &  & DCC &  \\ 
    & DSARH \cite{zhang2026dsarh} & & & DCC + PGD-based AT &  \\
    \cmidrule(lr){1-5}
    \multirow{10}{*}{\centering Continuous}
    & DCH \cite{caoDeepCauchyHashing2018} &\multirow{4}{*}{\centering $\min L_s(\bu) + h(\bu)$} & $\log \left( 1 + \frac{\text{dist}_{\text{h}}(|\bu|, \mathbf{1})}{\gamma} \right)$ & \multirow{4}{*}{\centering SGDM} &  \\
    & DNNH \cite{laiSimultaneousFeatureLearning2015} & & $\mathcal{I}_{[-1,1]^d}(\bu)$ & &  \\
    & DSH \cite{liuDeepSupervisedHashing} & & $\| | \bu| - \mathbf{1}\|_1$ &  &  \\ 
    & DTSH \cite{wang2016DTSH} & & $\| | \bu| - \mathbf{1}\|_2^2$ &  &  \\ 
    & DHN \cite{zhuDeepHashingNetwork2016} & & $\sum_i\log(\cosh(|u_i|-1))$&  &  \\ 
    \cmidrule(lr){2-5}
    & MDSHC \cite{wang2023deep} & \scriptsize $\min L_c(\bu) + L_s (\bu)+ h(\bu)$ & $\| | \bu| - \mathbf{1}\|_1$ & RMSProp &  \\ 
    \cmidrule(lr){2-5}
    & HashNet \cite{Cao2017HashNetDL} & \multirow{2}{*}{\centering $\min L_s(\bu)$} & \multirow{2}{*}{\centering -} & SGDM &  \\     
    & OrthoHash \cite{hoeOneLossAll2021} & & & Adam &  \\
    % \cmidrule(lr){2-5}
    % & SBPG \cite{wang2024bregman} & \scriptsize $\min L_s(\bu) + h(\bu)$ & Weakly convex $h$ & SBregman PG & $\mathcal{O}(\varepsilon^{-4})$ \\
    \cmidrule(lr){2-6}
   & \textbf{DualHash-StoM} & \multirow{2}{*}{\small $ \min L_s(\bu) + Q(\bu,\bb) + h(\bb)$} & \multirow{2}{*}{$\| | \bb| - \mathbf{1}\|_1$} & \multirow{2}{*}{Sto Primal-Dual} & $\mathcal{O}(\varepsilon^{-4})$ \\
    & \textbf{DualHash-StoRM} & & & & $\mathcal{O}(\varepsilon^{-3})$ \\
    \bottomrule
    \end{tabular}
    }
    % \vspace{-5pt}
    \parbox{\linewidth}{\scriptsize
    \textbf{Notes}: 
    The \textbf{Type} column indicates whether binary codes are explicitly updated; DSARH is grouped with Discrete because its unified-hashing branch uses DCC updates.
    The \textbf{Model} column shows the core optimization objective, where \(\bb\) denotes binary codes and \(\bu\) denotes continuous outputs. The \textbf{Quantization Loss} column specifies the loss (\(Q(\cdot, \cdot)\) or \(h(\cdot)\)) for quantization error, with `-' indicating no explicit loss. 
    % The \textbf{Optimization} column describes the training strategy, with discrete methods focusing on binary code updates. 
    DCC stands for Discrete Cyclic Coordinate, PGD-based AT for projected-gradient adversarial training, while BQP denotes Binary Quadratic Programming. The \textbf{Complexity} column reports the convergence rate, with `-' indicating no such analysis.
    }
     \vspace{-2ex}
\end{table}
\paragraph{Optimization in Deep Hashing}
% Existing approaches to bridging the continuous--discrete gap in deep hashing fall into three categories.
% One line of work adopts continuous relaxation and introduces explicit quantization regularization on network outputs; in particular, W-type regularizers such as $\||\bz|-\mathbf{1}\|_1$ have been widely used to encourage binary-like activations \cite{liuDeepSupervisedHashing,wang2023deep,zhuDeepHashingNetwork2016}. Another line refines the discretization objective through greedy coding schemes, central similarity targets, or proxy-based supervision \cite{su2018greedy,yuanCentralSimilarityQuantization2020,semanticHashCenters2025}. Alternative strategies directly address binary constraints by introducing auxiliary hash codes \cite{liDeepSupervisedHashing,li2017deep,liFeatureLearningBased2016}. Notably, even recent architectures retain W-type regularization; for instance, He and Wei \cite{he2025mambahash} employ $\||\bz|-\mathbf{1}\|_2^2$ as a smooth surrogate for $\||\bz|-\mathbf{1}\|_1$ to facilitate gradient-based optimization. These developments indicate that the practical relevance of W-type regularization has not diminished. The remaining challenge is to develop theoretically grounded optimization methods that can directly handle nonsmooth W-type regularization under nonlinear deep network mappings, which is precisely the focus of this work.
% Recent studies have further strengthened backbone architectures to improve feature quality, extending deep hashing from conventional CNNs to hybrid CNN--Transformer frameworks \cite{he2024hybridhash}, Transformer-based hashing models \cite{chen2021transhash,wu2025fth,liu2025capmh}, and state-space architectures \cite{he2025mambahash}. 
W-type regularizers such as \blue{$\bz\mapsto\lVert|\bz|-\mathbf{1}\rVert_1$} have been 
widely adopted in deep hashing to bridge the gap between continuous 
network outputs and discrete binary codes~\cite{liuDeepSupervisedHashing,
wang2023deep,zhuDeepHashingNetwork2016}.
Alternative strategies either introduce auxiliary hash codes with a 
smoother surrogate $\|\cdot\|_2^2$~\cite{liDeepSupervisedHashing,
li2017deep,liFeatureLearningBased2016}, or refine the discretization 
objective via greedy coding, central similarity targets, or proxy-based 
supervision~\cite{su2018greedy,yuanCentralSimilarityQuantization2020,
semanticHashCenters2025}. Nevertheless, even recent architectures continue to rely on W-type regularization to mitigate quantization error. For instance,
He and Wei~\cite{he2025mambahash} employ \blue{$\bz\mapsto\lVert|\bz|-\mathbf{1}\rVert_2^2$}
as a smooth surrogate for \blue{$\bz\mapsto\lVert|\bz|-\mathbf{1}\rVert_1$} to ease
gradient-based optimization, highlighting that the core difficulty of
developing theoretically grounded methods for nonsmooth W-type
regularization under nonlinear network mappings has not been resolved.
This is precisely the gap this work aims to fill.

\paragraph{Primal-Dual and Stochastic Splitting Methods} Primal-dual methods, such as PDHG~\cite{chambolleFirstOrderPrimalDualAlgorithm2011} and its unified extensions~\cite{esserZhangChan2010}, are foundational tools for composite optimization with broad applications in imaging sciences. These methods enjoy ergodic $\mathcal{O}(1/k)$ convergence and remain stable under adjoint mismatch~\cite{chouzenouxAdjointMismatch}. For convex nonlinear saddle-point models, Gao and Zhang~\cite{gao2023alternative} proposed an alternative extrapolation scheme with provable convergence guarantees. Stochastic counterparts, notably SPDHG~\cite{chambolleSPDHG2018} \blue{and stochastic majorization-minimization subspace methods~\cite{chouzenouxStochasticMM2017}}, handle structured finite-sum objectives with almost-sure and ergodic convergence. In the nonconvex setting, Clason et al.~\cite{clasonAccelerationGlobalConvergence2019} established global convergence under the Kurdyka--{\L}ojasiewicz (KL) condition, while Guo et al.~\cite{guo2023preconditioned} further addressed nonconvex stochastic composite problems. More recently, Shumaylov et al.~\cite{shumaylovweakly} incorporated weakly convex neural-network regularizers into a primal-dual framework and proved convergence of PDHGM to critical points under the KL condition. When the problem admits a block-separable structure, splitting methods such as ADMM with its stochastic variants~\cite{wangGlobalConvergenceADMM2018,ADMMNolinear2024,chen2025convergence,bian2021stochastic} and stochastic alternative minimization~\cite{driggs2021stochastic,InertialProximalAlternating,bian2026stochastic} provide promising approaches for nonconvex multi-block optimization.

\paragraph{Stochastic Nonconvex Constrained Optimization}
From a constrained optimization perspective, Problem~\eqref{Eq: the general nonlinear composite problem in intro} \blue{can be reformulated via a variable lifting technique~\cite{wang2025augmentedrnn,jin2025multicomposite} as a nonlinearly constrained optimization problem of the form
\begin{equation}
  \label{Eq: the constrained reformulation}
  \min_{\bw}\; f(\bw) + h(\bw)
  \quad \text{s.t.}\quad
  \bc(\bw) = 0.
\end{equation}}Several recent works have established oracle complexity results for first-order stochastic augmented-Lagrangian and penalty methods for problems of this type. For instance, Jin and Wang \cite{jin2022stochastic} proposed a single-loop stochastic primal-dual method based on a linearized augmented-Lagrangian function, achieving oracle complexities of $\mathcal{O}(\varepsilon^{-6})$ from an arbitrary initial point and $\mathcal{O}(\varepsilon^{-5})$ from a feasible one for locating an $\varepsilon$-KKT point. By incorporating the recursive momentum estimator of \cite{cutkoskyMomentumBasedVarianceReduction2019}, Shi et al. \cite{shi2025momentum} developed a momentum-based linearized augmented-Lagrangian method, improving these bounds to $\mathcal{O}(\varepsilon^{-4})$, or the nearly optimal $\mathcal{O}(\varepsilon^{-3})$ when initialized at a nearly feasible point. To certify a stronger notion of stationarity that deterministically enforces constraint violation of at most $\varepsilon$, Lu et al. \cite{lu2024variancereducedfirstordermethodsdeterministically} introduced a linearized quadratic-penalty method with truncated momentum. Under a regularity condition on the constraint gradients, recursive momentum yields the nearly optimal complexity $\widetilde{\mathcal{O}}(\varepsilon^{-3})$. 

\noindent\textbf{Notations.} We use $a$, $\ba$, and $\bA$ to denote scalars, vectors, and matrices, respectively.
\blue{All vectors are column vectors. For $\bB\in\RR^{n\times d}$, we write $\bB=[\bb_1^\top;\dots;\bb_n^\top]$, where $\bb_i\in\RR^d$ and $\bb_i^\top$ is the $i$-th row of $\bB$. The same convention applies to iterates: $\bB^k=[(\bb_1^k)^\top;\dots;(\bb_n^k)^\top]$.}
Throughout this paper, \blue{$\bm{1}$ denotes an all-ones vector or matrix of the appropriate size}, and $i, j$ serve as indices. For matrix variables \(\bB\), \(\bY\), and \(\bLmbd\), we work directly in \(\RR^{n \times d}\), equipped with the Frobenius inner product \(\langle \cdot, \cdot \rangle\) and induced Frobenius norm \(\| \cdot \| := \sqrt{\langle \cdot, \cdot \rangle}\). Unless specified otherwise, \(\| \cdot \|\) denotes the Euclidean norm for vectors and the Frobenius norm for matrices. \( |\cdot| \) denotes element-wise absolute value. \blue{For a matrix $\bA=[A_{ij}]$, $\|\bA\|_1:=\sum_{i,j}|A_{ij}|$ denotes the entrywise $\ell_1$ norm.} For a closed set \(\calC \subseteq \RR^{q}\), the distance from a point \(\bx \in \RR^{q}\) to \(\calC\) is defined as \(\text{dist}(\bx, \mathcal{C}) := \min_{\by \in \mathcal{C}}\|\bx - \by\|\).
\blue{For simplicity, we use the same symbols $h$ and $h^*$ for their scalar forms and their separable extensions to vectors and matrices whenever no confusion can arise. The Fenchel conjugate of $h$~\cite{Rockafellar1998VariationalA} is $h^*(\bLmbd) := \sup_{\bZ \in \RR^{n\times d}} \{ \langle \bLmbd, \bZ \rangle - h(\bZ) \}$ for $\bLmbd \in \RR^{n\times d}$.}
